\documentclass[12pt,twoside,a4paper,titlepage]{article}
\usepackage[utf8]{inputenc}
\usepackage[T1]{fontenc}
\usepackage[english,francais]{babel}
\usepackage{listings} % introduction de code
\usepackage{geometry}
\usepackage[colorlinks=true,linkcolor=black,citecolor=blue,urlcolor=blue]{hyperref} %URLs,index etc etc
\usepackage{amsmath} %Formules mathématiques
\usepackage{textcomp}
\usepackage{graphicx}
\usepackage{caption}
\usepackage{fullpage}
\usepackage{fancyhdr}  %entêtes et pieds de page
\usepackage{xcolor}
\usepackage{tikz} %Schémas
\usetikzlibrary{arrows}

\setcounter{tocdepth}{1} % On veut uniquement les sections et pas les sous sections dans la table des matières
\setlength{\doublerulesep}{\arrayrulewidth} % Lignes épaisses dans les tableaux

\begin{document}

%%%%%%%%%%%%%%%%%%%%%%%%%%%%%%%%%%%%%%%
%%%%   Définitions header&footer   %%%%
%%%%%%%%%%%%%%%%%%%%%%%%%%%%%%%%%%%%%%%
\pagestyle{fancy}
\fancyhf{} % on efface tout et on recommence
% PIED DE PAGE :
\fancyfoot[LO,RE]{Groupe 52 \hfill{} \textbf{Projet Informatique
								2\up{ème} Année \hfill{} \thepage }}
% épaisseur des traits
\renewcommand{\headrulewidth}{0pt}
\renewcommand{\footrulewidth}{2pt}

% le cas de la première page d’un chapitre
\fancypagestyle{plain}{%
	\fancyhf{}% on efface tout
	\fancyfoot[LO,RE]{Groupe 52 \hfill{} \textbf{Projet Informatique
						2\up{ème} Année} \hfill{} \thepage }
	% on efface tous les traits
	\renewcommand{\headrulewidth}{0pt}%
	\renewcommand{\footrulewidth}{2pt}
}

%%%%%%%%%%%%%%%%%%%%%%%%%%%%%%%%%%%%%%%%%%%%%%%%%%%%%%%%%%%%%
%%%%%%%%          Page de Garde          %%%%%%%%%%%%%%%%%%%%
%%%%%%%%%%%%%%%%%%%%%%%%%%%%%%%%%%%%%%%%%%%%%%%%%%%%%%%%%%%%%
\begin{minipage}[c]{0.5\linewidth}
\includegraphics[bb=0 50 200 100, scale=0.3]{../images/insa.JPG}
\end{minipage}
\begin{minipage}[r]{0.5\linewidth}
\vspace*{25pt}
\begin{flushright}
\LARGE D\'epartement \\ du Premier \\ Cycle \\ \rule{\linewidth}{2pt}
\end{flushright}
\end{minipage}

\vspace*{1cm}
\thispagestyle{empty}
\begin{center}
\LARGE Projet Informatique 2\up{\`eme} Ann\'ee\\
\end{center}
\vspace*{1.3cm}
	
\setlength\fboxrule{2pt}
\fbox{
\begin{minipage}{0.95\textwidth}
     \begin{center}
     \vspace*{0.21cm}\LARGE \color{blue} \bsc{Simulation mécanique d'un solide\\
     et\\
     moteur graphique 2D en Pascal}
     \end{center}
\end{minipage}
}


\begin{center}
\vspace*{1cm}
%Changer les noms ici
\LARGE \color{black} Tristan \bsc{Pourcelot} \\
					 Victor \bsc{Cleren} \\
					 Rémi \bsc{Verschelde} \\
					 Th\'eo \bsc{Garcia-Guitton} \\
\vspace*{1.5cm}
\Large Responsable : C. \bsc{Wolf}\\

\vspace*{1cm}

\Large Groupe 52 \hfill Ann\'ee scolaire 2010-2011\\
\vspace*{1cm}
\textbf{\Large Institut National des Sciences Appliqu\'{e}es de Lyon} \pagebreak
\end{center}

%%%%%%%%%%%%%%%%%%%%%%%%%%%%%%%%%%%%%%%
%%%%	   Corps du doc 		   %%%%
%%%%%%%%%%%%%%%%%%%%%%%%%%%%%%%%%%%%%%%
\section{Introduction}
Nous sommes partis d'une idée simple, celle de modéliser la chute d'un objet sur un
décor, sachant que l'utilisateur pouvait influer à la fois sur le décor et sur l'objet.
Nous avons également voulu intégrer de nombreux paramètres tels que la 
variation du matériau constituant le décor et l'objet, ainsi que la forme de cet objet.
L'utilisateur pourrait ainsi (au prix de quelques approximations)
simuler la chute d'un satellite sur la Lune.

\section{Spécifications}
L'objet principal de notre programme est de modéliser la chute d'un solide dans un
fluide, tout en prenant en compte son interaction avec un environnement solide (chocs).

\subsection{IHM - \emph{Interface Homme - Machine}}
Il est proposé à l'utilisateur de dessiner à la fois le décor et l'objet.
Les deux dessins sont séparés en deux fenêtres, munie chacunes de leur interface propre,
puisque les outils nécessaires au dessin du décor ne sont pas les mêmes que pour le dessin de l'objet.
L'utilisateur peut également ouvrir une fenêtre dédiée aux paramètres physiques, afin d'influer
sur la valeur de la gravité notamment.

\subsection{Problèmes physiques et méthodes de résolution}
\subsubsection{Forces}
Le programme doit prendre en compte les différentes forces appliquées à 
un instant $t$ sur le solide. Le mouvement de ce solide est donc déterminé
par les équations de la mécanique newtonienne.
On peut donc définir les différentes forces appliquées sur le solide :

\begin{itemize}
\item[$\bullet$] La gravité : \dotfill{} $\overrightarrow{P} = mg\, \overrightarrow{y_0}$ 
\item[$\bullet$] La poussée d'Archimède : \dotfill{} $\overrightarrow{F_{Ar}} = -\rho_{liquide}\cdot V_{deplace} \, \overrightarrow{y_0} $
\item[$\bullet$] Les frottements lors du contact: \dotfill{}$\overrightarrow{F_{fr}} = -f \,\overrightarrow{V}_{\text{Objet/Décor}}$
\end{itemize}

La majorité des forces sont résolues à l'aide d'algorithmes basiques de
calcul, prenant en compte des paramètres telles que le volume de l'objet,
sa vitesse ou son poids.
\subsubsection{\emph{PFD} - Principe Fondamental de la Dynamique }
Dans le cas d'un problème plan, notre problème s'exprime sous la forme :
\begin{equation}
\sum F_x = m \frac {dv_x}{dt} \: \text{et} \: \sum F_y = m\frac{dv_y}{dt} \: \text{et} \: \sum M_z = J \frac {d\omega}{dt}
\end{equation}
En supposant le $\Delta t$ du timer suffisament faible, on peut assimiler
$\frac {dv_x}{dt}$ à $\frac{\Delta v_x}{\Delta t}$,
$ \frac {dv_y}{dt} \approx \frac{\Delta v_y}{\Delta t}$ et
$ \frac{d\omega}{dt} \approx \frac{\Delta \omega}{\Delta t}$ .
\newline
Avec les mêmes approximations, on obtient :
$\frac {dx}{dt} = v_x \approx \frac{\Delta x}{\Delta t}$, 
$ \frac {dy}{dt} \approx \frac{\Delta y}{\Delta t}$ et
$ \frac{d\theta}{dt} \approx \frac{\Delta \theta}{\Delta t}$. \\
La vitesse initiale $\overrightarrow{v_0}$ est définie par l'utilisateur.
Ainsi, connaissant la vitesse $\overrightarrow{v(t)}$ à un instant $t$,
on peut calculer la vitesse $\overrightarrow{v(t+\Delta t)}$ à l'instant $t + \Delta t$.

\begin{equation}
v_{x}(t+\Delta t)= v_{x}(t) + m \Delta t \sum F_x \: \text{et}\:
v_{y}(t+\Delta t)= v_{y}(t) + m\Delta t\sum F_y \: \text{et}\:
\omega (t + \Delta t)= \omega(t) + J\Delta t \sum m_z 
\end{equation}
En prenant en compte les approximations précédentes, on obtient le système suivant :
\begin{equation}
x(t+\Delta t) = x(t) +\Delta t \cdot v_{x}(t) \: \text{et} \:
y(t+\Delta t) = x(t) +\Delta t \cdot v_{x}(t) \: \text{et} \:
\theta (t+\Delta t) = \theta (t) + \Delta t \cdot \omega (t) \: \text \:
\end{equation}

\newpage
\subsubsection{Collisions}
Notons $\overrightarrow{v_c}$ la vitesse du point de contact et $\overrightarrow{v_g}$
la vitesse, connue, du centre de gravité du solide.
Pour le calcul de la nouvelle vitesse du solide, nous supposons que le contact
entre le solide et le décor lors du choc est ponctuel. 
On se place dans la base de contact définie par le schéma suivant :

\begin{center}
\begin{figure}[h]
\begin{center}
\scalebox{0.2}[0.2]{\includegraphics*{../images/collision.png}}
\end{center}
\caption{Base locale du contact}
\end{figure}
\end{center}
Le contact étant ponctuel, on suppose que les forces appliquées au point
de contact C sont similaires à celles appliquées dans le modèle de choc d'une particule, on a donc :
\begin{equation}
v_{C,x,initial} = v_{C,x,final} \: \text{et} \: v_{C,y,initial} = -v_{C,y,final}
\end{equation}
Or, en notant $x_G$ et $y_G$ les coordonnées du centre de gravité dans 
la base de contact :
\begin{equation}
v_{C,x}=v_{G,x}+y_G \cdot \omega \: \text{et} \: v_{C,y} = v_{g,y} -x_G\cdot \omega
\end{equation}
On suppose dans un premier temps que l'énergie cinétique est conservée :
\begin{equation}
J\omega^2_{initial} + m(v^2_{g,x,initial} + v^2_{g,y,initial} )
=
J\omega^2_{final} + m(v^2_{g,x,final} + v^2_{g,y,final} )
\end{equation}
On obtient ainsi un système de 3 équations qui nous donne les équations de
mouvement après collision.
$(v_{x0} \: , \: v_{x1})$ désigne $(v_{G,x,initial} \: , \: v_{G,y,initial})$,
de même pour $y$ et $\omega$ .\\
On obtient ainsi les solutions suivantes :
\begin{equation}
\omega_1 = \frac{-m(x^2_G+v_{x0}y_G+y^2_G\omega_0 + v_{y0}) +\sqrt{\Delta}}{my^2_G + mx^2_G +J}
\end{equation}

\begin{eqnarray*}
	\Delta & = & \omega^2 J^2 -2Jmv_{x0}y_G\omega_0 + 2Jmv_{y0}x_G \omega_0  + 2m^2v_{x0}y_Gv_{y0}x_G - 4m^2v_{x0}y_Gx^2_G\omega_0 \\
		& & +\ 4m^2y^2_g\omega_0 v_{y0} x_G- 4m^2y^2_g\omega^2_0x^2_g + m^2v^2_{x0}y^2_G + m^2v^2_{y0}x^2_G
\end{eqnarray*}

On obtient ainsi :
\begin{equation}
v_{x1}=v_{x0} +y_G(\omega_0 - \omega_1) \: \ \text{et} \ \:  v_{y1}=v_{y0} +x_G(\omega_0 - \omega_1)
\end{equation} 

\newpage

\section{Réalisation}
\subsection{Interface}
La fenêtre principale affiche un menu utilisateur, une fenêtre de dessin,
et un espace d'affichage des principaux paramètres physiques.
Outre ces fonctionnalités, c'est également à elle d'effectuer l'animation
de l'image, en calculant régulièrement (présence d'un Timer) la position
du solide ainsi que sa vitesse et son interaction avec le décor.
Les autres fenêtres ont pour mission de gérer le dessin de l'objet et son remplissage,
le dessin du décor et la gestion des matériaux.


\subsection{Gestion des chocs}
La partie la plus compliquée dans la gestion des chocs reste la détection des chocs
et le calcul de la tangente locale au décor.

A chaque itération du timer, une procédure est lancée après le calcul de 
la nouvelle position de l'objet pour détecter une collision potentielle.
On parcourt la matrice du solide et l'extrait correspondant de la matrice
du décor pour chercher d'éventuelles superpositions de pixels non blancs.
Les dimensions horizontales et verticales de l'objet ont été au préalable
calculées pour optimiser ce parcours.

Il s'agit ensuite de déterminer la zone restreinte du contact pour la réduire ensuite
à un unique pixel qui servira dans les calculs dynamiques.
Pour ce faire, une boucle fait reculer (sans l'afficher) le solide dans 
la direction opposée à la vitesse, pixel par pixel,
jusqu'à ce qu'il n'y ait plus de collision.

La procédure fait aussi tourner l'objet petit à petit pour retrouver
son angle de rotation avant le choc.
L'ensemble des pixels colorés du décor superposés aux points du solide
sont stockés dans un tableau. Lorsque la boucle se termine, ce tableau
contient les points de la surface du décor qui délimitent la zone de contact.
On prend ensuite un point au milieu de cette zone que l'on définit comme pixel de collision.\\

\begin{center}
\begin{figure}[h]
\begin{center}
\scalebox{0.3}[0.3]{\includegraphics*{../images/tangente.png}}
\end{center}
\caption{Schématisation du calcul de la tangente}
\end{figure}
\end{center}

Le calcul des tangentes se fait à partir du pixel de contact et des deux pixels
constituant la surface les plus proches du pixel de contact.

Pour déterminer ces deux pixels, la procédure teste chaque pixel adjacent
au pixel de contact et s'arrête sur les premiers pixels pleins trouvés.

Les deux premiers pixels testés sont déterminés par la vitesse à laquelle
l'objet arrive vers la surface (cf schéma ci-dessus).

Si aucun de ces pixels n'est plein (ce qui en pratique devrait toujours être le cas),
une procédure teste tour à tour les pixels en tournant autour du point de contact,
dans le sens horaire pour un point et dans le sens trigonométrique pour l'autre.

Une fois les deux points déterminés, le vecteur tangent est calculé en
faisant la moyenne du vecteur reliant le point de contact et le premier point
et du vecteur reliant le point de contact et le deuxième point.



\subsection{Remplissage de l'objet}
Un des problèmes majeurs qui nous fut posé durant ce projet fut le remplissage
du solide afin de pouvoir déterminer sa masse et son centre d'inertie.
Tout d'abord, il fallut s'assurer que le dessin de l'utilisateur était bien connexe.
Pour remédier à ce problème, nous traçons une droite entre le point de départ
$(X_e,Y_e)$ et le point de sortie $(X_s,Y_s)$.
Nous avons choisi la méthode du ScanLine-FloodFill pour remplir le solide dessiné.
Cette méthode utilise une graine (la \emph{seed}) placée sur un point quelquonque
de notre solide, et analyse tout les pixels environnants pour déterminer
quels sont ceux à l'intérieur du solide, et les autres.
Après plusieurs essais infructueux avec un algorithme classique, nous avons
adapté un algorithme en C trouvé sur Internet, qui analyse ligne par ligne l'objet,
permettant ainsi un gain considérable en efficacité. 
%
%\begin{algorithm}[h]
%\caption{Pseudo algorithme de remplissage par ligne}
%\begin{algorithmic}
%	\STATE Prout;
%	\PROCEDURE {RemplirLigne} {$(aOldColor,aNewColor, X_{seed}, Y_{seed})$}
%		\WHILE{$couleur(Y+1) = aOldColor$ and $couleur(Y+1) \neq couleur_{\text{frontière}}$}
%			\STATE $couleur(Y+1) \gets aNewColor$
%		\ENDWHILE
%		\IF{$couleur_{voisin\_ gauche} = aOldColor$}
%			\STATE remplir\_ Ligne(X-1,Y)
%		\ENDIF
%		\IF{$couleur_{voisin\_ droit} = aOldColor$}
%			\STATE remplir\_ Ligne(X+1,Y)
%		\ENDIF
%	\ENDPROCEDURE
%\end{algorithmic}
%\end{algorithm}
%\begin{algorithmic}[1]
%\repeat
%\Comment{forever}
%\State this\Until{you die.}
%\end{algorithmic}

\newpage

\section{Bugs - Améliorations}
Il subsiste de nombreux bugs et améliorations possibles de notre projet.
En particulier, de nombreuses améliorations que nous pensions apporter
à notre projet n'ont pas été implémentées par faute de temps, certains 
problèmes ayant été plus chronophages que prévu.
\subsection{Améliorations}
\begin{itemize}
\item[$\bullet$] Différents fluides (air, eau, vide \dots)
\item[$\bullet$] Variation aléatoire et incidence du vent 
\item[$\bullet$] Amélioration du modèle physique (élasticité \dots )
\item[$\bullet$] Modification du décor et de l'objet en temps réel par l'utilisateur
\item[$\bullet$] Tracé de la trajectoire de l'objet
\item[$\bullet$] Force de frottement prenant en compte la forme de l'objet
\item[$\bullet$] Roulement ou arrêt de l'objet possible
\item[$\bullet$] Possibilité de modification des paramètres physiques (prévu mais pas implémenté)
\end{itemize}
\subsection{Bugs}
\begin{itemize}
\item[$\bullet$] Lenteur de l'exécution sous Linux (pas de problèmes sous Windows)
\item[$\bullet$] Le décor ne se remet pas à jour si on le redessine
\item[$\bullet$] Calcul des tangentes peu précis (sinon il serait beaucoup trop lent)
\item[$\bullet$] Mauvais fonctionnement sur certains ordinateurs (64 bits?)
\end{itemize}
 
\section{Conclusion}
All your base are belong to us!
\newpage

\section{Annexes}
\subsection{Avis personnels et problèmes rencontrés}
\paragraph{Tristan :}
\begin{itemize}
\item[$\bullet$]Communication (qui fait quoi, quand et pourquoi\dots )
\item[$\bullet$]Conflits SVN
\item[$\bullet$]Implémentation pascal linux = caca
\end{itemize}
\paragraph{Rémi :}
\begin{itemize}
\item[$\bullet$] Les autres ont rien foutus
\end{itemize}
\paragraph{Victor :}
\begin{itemize}
\item Idem
\end{itemize}

\paragraph{Théo :}
<<<<<<< .mine
\begin{itemize}
\item niah
=======
\begin{itemize}
\item niah
\itemindent Prout
\end{itemize}
>>>>>>> .r138

<<<<<<< .mine
\end{itemize}

=======
>>>>>>> .r138
\newpage
\subsection{Diagrammes de classes}

\begin{center}
\begin{figure}[h]
\begin{center}
\scalebox{0.5}[0.42]{\includegraphics*{../images/UML.jpeg}}
\end{center}
\caption{Diagramme de classe utilisé pour notre projet}
\end{figure}
\end{center}
%% Graphic for TeX using PGF
% Title: /home/blackout/Documents/Sources/2010_pascal_moteur_2D/trunk/doc/Diagramme UML.dia
% Creator: Dia v0.97.1
% CreationDate: Sat May 28 23:45:59 2011
% For: blackout
% \usepackage{tikz}
% The following commands are not supported in PSTricks at present
% We define them conditionally, so when they are implemented,
% this pgf file will use them.
\ifx\du\undefined
  \newlength{\du}
\fi
\setlength{\du}{15\unitlength}
\begin{tikzpicture}
\pgftransformxscale{1.000000}
\pgftransformyscale{-1.000000}
\definecolor{dialinecolor}{rgb}{0.000000, 0.000000, 0.000000}
\pgfsetstrokecolor{dialinecolor}
\definecolor{dialinecolor}{rgb}{1.000000, 1.000000, 1.000000}
\pgfsetfillcolor{dialinecolor}
\pgfsetlinewidth{0.100000\du}
\pgfsetdash{}{0pt}
\definecolor{dialinecolor}{rgb}{1.000000, 1.000000, 1.000000}
\pgfsetfillcolor{dialinecolor}
\fill (31.550000\du,-8.750000\du)--(31.550000\du,-7.350000\du)--(42.445000\du,-7.350000\du)--(42.445000\du,-8.750000\du)--cycle;
\definecolor{dialinecolor}{rgb}{0.000000, 0.000000, 0.000000}
\pgfsetstrokecolor{dialinecolor}
\draw (31.550000\du,-8.750000\du)--(31.550000\du,-7.350000\du)--(42.445000\du,-7.350000\du)--(42.445000\du,-8.750000\du)--cycle;
% setfont left to latex
\definecolor{dialinecolor}{rgb}{0.000000, 0.000000, 0.000000}
\pgfsetstrokecolor{dialinecolor}
\node at (36.997500\du,-7.800000\du){CTorseur};
\definecolor{dialinecolor}{rgb}{1.000000, 1.000000, 1.000000}
\pgfsetfillcolor{dialinecolor}
\fill (31.550000\du,-7.350000\du)--(31.550000\du,2.600000\du)--(42.445000\du,2.600000\du)--(42.445000\du,-7.350000\du)--cycle;
\definecolor{dialinecolor}{rgb}{0.000000, 0.000000, 0.000000}
\pgfsetstrokecolor{dialinecolor}
\draw (31.550000\du,-7.350000\du)--(31.550000\du,2.600000\du)--(42.445000\du,2.600000\du)--(42.445000\du,-7.350000\du)--cycle;
% setfont left to latex
\definecolor{dialinecolor}{rgb}{0.000000, 0.000000, 0.000000}
\pgfsetstrokecolor{dialinecolor}
\node[anchor=west] at (31.700000\du,-6.650000\du){\#fX: real};
% setfont left to latex
\definecolor{dialinecolor}{rgb}{0.000000, 0.000000, 0.000000}
\pgfsetstrokecolor{dialinecolor}
\node[anchor=west] at (31.700000\du,-5.925000\du){Composante};
\definecolor{dialinecolor}{rgb}{0.000000, 0.000000, 0.000000}
\pgfsetstrokecolor{dialinecolor}
\node[anchor=west] at (31.700000\du,-5.225000\du){horizontale de};
\definecolor{dialinecolor}{rgb}{0.000000, 0.000000, 0.000000}
\pgfsetstrokecolor{dialinecolor}
\node[anchor=west] at (31.700000\du,-4.525000\du){l'effort};
% setfont left to latex
\definecolor{dialinecolor}{rgb}{0.000000, 0.000000, 0.000000}
\pgfsetstrokecolor{dialinecolor}
\node[anchor=west] at (31.700000\du,-3.400000\du){+fY: real};
% setfont left to latex
\definecolor{dialinecolor}{rgb}{0.000000, 0.000000, 0.000000}
\pgfsetstrokecolor{dialinecolor}
\node[anchor=west] at (31.700000\du,-2.675000\du){Composante};
\definecolor{dialinecolor}{rgb}{0.000000, 0.000000, 0.000000}
\pgfsetstrokecolor{dialinecolor}
\node[anchor=west] at (31.700000\du,-1.975000\du){verticale de};
\definecolor{dialinecolor}{rgb}{0.000000, 0.000000, 0.000000}
\pgfsetstrokecolor{dialinecolor}
\node[anchor=west] at (31.700000\du,-1.275000\du){l'effort};
% setfont left to latex
\definecolor{dialinecolor}{rgb}{0.000000, 0.000000, 0.000000}
\pgfsetstrokecolor{dialinecolor}
\node[anchor=west] at (31.700000\du,-0.150000\du){+fMz: real};
% setfont left to latex
\definecolor{dialinecolor}{rgb}{0.000000, 0.000000, 0.000000}
\pgfsetstrokecolor{dialinecolor}
\node[anchor=west] at (31.700000\du,0.575000\du){Composante du};
\definecolor{dialinecolor}{rgb}{0.000000, 0.000000, 0.000000}
\pgfsetstrokecolor{dialinecolor}
\node[anchor=west] at (31.700000\du,1.275000\du){moment normale};
\definecolor{dialinecolor}{rgb}{0.000000, 0.000000, 0.000000}
\pgfsetstrokecolor{dialinecolor}
\node[anchor=west] at (31.700000\du,1.975000\du){au plan};
\definecolor{dialinecolor}{rgb}{1.000000, 1.000000, 1.000000}
\pgfsetfillcolor{dialinecolor}
\fill (31.550000\du,2.600000\du)--(31.550000\du,3.600000\du)--(42.445000\du,3.600000\du)--(42.445000\du,2.600000\du)--cycle;
\definecolor{dialinecolor}{rgb}{0.000000, 0.000000, 0.000000}
\pgfsetstrokecolor{dialinecolor}
\draw (31.550000\du,2.600000\du)--(31.550000\du,3.600000\du)--(42.445000\du,3.600000\du)--(42.445000\du,2.600000\du)--cycle;
% setfont left to latex
\definecolor{dialinecolor}{rgb}{0.000000, 0.000000, 0.000000}
\pgfsetstrokecolor{dialinecolor}
\node[anchor=west] at (31.700000\du,3.300000\du){+calculForce(aForme:CForme)};
\pgfsetlinewidth{0.100000\du}
\pgfsetdash{}{0pt}
\definecolor{dialinecolor}{rgb}{1.000000, 1.000000, 1.000000}
\pgfsetfillcolor{dialinecolor}
\fill (10.600000\du,4.050000\du)--(10.600000\du,6.850000\du)--(26.115000\du,6.850000\du)--(26.115000\du,4.050000\du)--cycle;
\definecolor{dialinecolor}{rgb}{0.000000, 0.000000, 0.000000}
\pgfsetstrokecolor{dialinecolor}
\draw (10.600000\du,4.050000\du)--(10.600000\du,6.850000\du)--(26.115000\du,6.850000\du)--(26.115000\du,4.050000\du)--cycle;
% setfont left to latex
\definecolor{dialinecolor}{rgb}{0.000000, 0.000000, 0.000000}
\pgfsetstrokecolor{dialinecolor}
\node at (18.357500\du,5.000000\du){CPoids};
% setfont left to latex
\definecolor{dialinecolor}{rgb}{0.000000, 0.000000, 0.000000}
\pgfsetstrokecolor{dialinecolor}
\node at (18.357500\du,5.775000\du){Effet lié au};
\definecolor{dialinecolor}{rgb}{0.000000, 0.000000, 0.000000}
\pgfsetstrokecolor{dialinecolor}
\node at (18.357500\du,6.475000\du){poids};
\definecolor{dialinecolor}{rgb}{1.000000, 1.000000, 1.000000}
\pgfsetfillcolor{dialinecolor}
\fill (10.600000\du,6.850000\du)--(10.600000\du,9.600000\du)--(26.115000\du,9.600000\du)--(26.115000\du,6.850000\du)--cycle;
\definecolor{dialinecolor}{rgb}{0.000000, 0.000000, 0.000000}
\pgfsetstrokecolor{dialinecolor}
\draw (10.600000\du,6.850000\du)--(10.600000\du,9.600000\du)--(26.115000\du,9.600000\du)--(26.115000\du,6.850000\du)--cycle;
% setfont left to latex
\definecolor{dialinecolor}{rgb}{0.000000, 0.000000, 0.000000}
\pgfsetstrokecolor{dialinecolor}
\node[anchor=west] at (10.750000\du,7.550000\du){\#fg: real};
% setfont left to latex
\definecolor{dialinecolor}{rgb}{0.000000, 0.000000, 0.000000}
\pgfsetstrokecolor{dialinecolor}
\node[anchor=west] at (10.750000\du,8.275000\du){Constante};
\definecolor{dialinecolor}{rgb}{0.000000, 0.000000, 0.000000}
\pgfsetstrokecolor{dialinecolor}
\node[anchor=west] at (10.750000\du,8.975000\du){gravitationnelle};
\definecolor{dialinecolor}{rgb}{1.000000, 1.000000, 1.000000}
\pgfsetfillcolor{dialinecolor}
\fill (10.600000\du,9.600000\du)--(10.600000\du,10.600000\du)--(26.115000\du,10.600000\du)--(26.115000\du,9.600000\du)--cycle;
\definecolor{dialinecolor}{rgb}{0.000000, 0.000000, 0.000000}
\pgfsetstrokecolor{dialinecolor}
\draw (10.600000\du,9.600000\du)--(10.600000\du,10.600000\du)--(26.115000\du,10.600000\du)--(26.115000\du,9.600000\du)--cycle;
% setfont left to latex
\definecolor{dialinecolor}{rgb}{0.000000, 0.000000, 0.000000}
\pgfsetstrokecolor{dialinecolor}
\node[anchor=west] at (10.750000\du,10.300000\du){+calculForce(aForme:CForme,aMasse:real)};
\pgfsetlinewidth{0.100000\du}
\pgfsetdash{}{0pt}
\definecolor{dialinecolor}{rgb}{1.000000, 1.000000, 1.000000}
\pgfsetfillcolor{dialinecolor}
\fill (19.400000\du,19.100000\du)--(19.400000\du,22.600000\du)--(37.225000\du,22.600000\du)--(37.225000\du,19.100000\du)--cycle;
\definecolor{dialinecolor}{rgb}{0.000000, 0.000000, 0.000000}
\pgfsetstrokecolor{dialinecolor}
\draw (19.400000\du,19.100000\du)--(19.400000\du,22.600000\du)--(37.225000\du,22.600000\du)--(37.225000\du,19.100000\du)--cycle;
% setfont left to latex
\definecolor{dialinecolor}{rgb}{0.000000, 0.000000, 0.000000}
\pgfsetstrokecolor{dialinecolor}
\node at (28.312500\du,20.050000\du){CFrottement};
% setfont left to latex
\definecolor{dialinecolor}{rgb}{0.000000, 0.000000, 0.000000}
\pgfsetstrokecolor{dialinecolor}
\node at (28.312500\du,20.825000\du){Frottements liés};
\definecolor{dialinecolor}{rgb}{0.000000, 0.000000, 0.000000}
\pgfsetstrokecolor{dialinecolor}
\node at (28.312500\du,21.525000\du){au fluide};
\definecolor{dialinecolor}{rgb}{0.000000, 0.000000, 0.000000}
\pgfsetstrokecolor{dialinecolor}
\node at (28.312500\du,22.225000\du){environnant};
\definecolor{dialinecolor}{rgb}{1.000000, 1.000000, 1.000000}
\pgfsetfillcolor{dialinecolor}
\fill (19.400000\du,22.600000\du)--(19.400000\du,25.350000\du)--(37.225000\du,25.350000\du)--(37.225000\du,22.600000\du)--cycle;
\definecolor{dialinecolor}{rgb}{0.000000, 0.000000, 0.000000}
\pgfsetstrokecolor{dialinecolor}
\draw (19.400000\du,22.600000\du)--(19.400000\du,25.350000\du)--(37.225000\du,25.350000\du)--(37.225000\du,22.600000\du)--cycle;
% setfont left to latex
\definecolor{dialinecolor}{rgb}{0.000000, 0.000000, 0.000000}
\pgfsetstrokecolor{dialinecolor}
\node[anchor=west] at (19.550000\du,23.300000\du){\#fCoef: real};
% setfont left to latex
\definecolor{dialinecolor}{rgb}{0.000000, 0.000000, 0.000000}
\pgfsetstrokecolor{dialinecolor}
\node[anchor=west] at (19.550000\du,24.025000\du){Coefficient de};
\definecolor{dialinecolor}{rgb}{0.000000, 0.000000, 0.000000}
\pgfsetstrokecolor{dialinecolor}
\node[anchor=west] at (19.550000\du,24.725000\du){frottement};
\definecolor{dialinecolor}{rgb}{1.000000, 1.000000, 1.000000}
\pgfsetfillcolor{dialinecolor}
\fill (19.400000\du,25.350000\du)--(19.400000\du,27.150000\du)--(37.225000\du,27.150000\du)--(37.225000\du,25.350000\du)--cycle;
\definecolor{dialinecolor}{rgb}{0.000000, 0.000000, 0.000000}
\pgfsetstrokecolor{dialinecolor}
\draw (19.400000\du,25.350000\du)--(19.400000\du,27.150000\du)--(37.225000\du,27.150000\du)--(37.225000\du,25.350000\du)--cycle;
% setfont left to latex
\definecolor{dialinecolor}{rgb}{0.000000, 0.000000, 0.000000}
\pgfsetstrokecolor{dialinecolor}
\node[anchor=west] at (19.550000\du,26.050000\du){+calculForce(aForme:CForme,aVitesse:CVitesse)};
% setfont left to latex
\definecolor{dialinecolor}{rgb}{0.000000, 0.000000, 0.000000}
\pgfsetstrokecolor{dialinecolor}
\node[anchor=west] at (19.550000\du,26.850000\du){+calculCoef(aForme:CForme)};
\pgfsetlinewidth{0.100000\du}
\pgfsetdash{}{0pt}
\definecolor{dialinecolor}{rgb}{1.000000, 1.000000, 1.000000}
\pgfsetfillcolor{dialinecolor}
\fill (47.300000\du,4.000000\du)--(47.300000\du,7.500000\du)--(58.195000\du,7.500000\du)--(58.195000\du,4.000000\du)--cycle;
\definecolor{dialinecolor}{rgb}{0.000000, 0.000000, 0.000000}
\pgfsetstrokecolor{dialinecolor}
\draw (47.300000\du,4.000000\du)--(47.300000\du,7.500000\du)--(58.195000\du,7.500000\du)--(58.195000\du,4.000000\du)--cycle;
% setfont left to latex
\definecolor{dialinecolor}{rgb}{0.000000, 0.000000, 0.000000}
\pgfsetstrokecolor{dialinecolor}
\node at (52.747500\du,4.950000\du){CArchimede};
% setfont left to latex
\definecolor{dialinecolor}{rgb}{0.000000, 0.000000, 0.000000}
\pgfsetstrokecolor{dialinecolor}
\node at (52.747500\du,5.725000\du){Effort lié à la};
\definecolor{dialinecolor}{rgb}{0.000000, 0.000000, 0.000000}
\pgfsetstrokecolor{dialinecolor}
\node at (52.747500\du,6.425000\du){poussée};
\definecolor{dialinecolor}{rgb}{0.000000, 0.000000, 0.000000}
\pgfsetstrokecolor{dialinecolor}
\node at (52.747500\du,7.125000\du){d'Archimède};
\definecolor{dialinecolor}{rgb}{1.000000, 1.000000, 1.000000}
\pgfsetfillcolor{dialinecolor}
\fill (47.300000\du,7.500000\du)--(47.300000\du,10.950000\du)--(58.195000\du,10.950000\du)--(58.195000\du,7.500000\du)--cycle;
\definecolor{dialinecolor}{rgb}{0.000000, 0.000000, 0.000000}
\pgfsetstrokecolor{dialinecolor}
\draw (47.300000\du,7.500000\du)--(47.300000\du,10.950000\du)--(58.195000\du,10.950000\du)--(58.195000\du,7.500000\du)--cycle;
% setfont left to latex
\definecolor{dialinecolor}{rgb}{0.000000, 0.000000, 0.000000}
\pgfsetstrokecolor{dialinecolor}
\node[anchor=west] at (47.450000\du,8.200000\du){\#fRho: real};
% setfont left to latex
\definecolor{dialinecolor}{rgb}{0.000000, 0.000000, 0.000000}
\pgfsetstrokecolor{dialinecolor}
\node[anchor=west] at (47.450000\du,8.925000\du){Masse volumique};
\definecolor{dialinecolor}{rgb}{0.000000, 0.000000, 0.000000}
\pgfsetstrokecolor{dialinecolor}
\node[anchor=west] at (47.450000\du,9.625000\du){du fluide};
\definecolor{dialinecolor}{rgb}{0.000000, 0.000000, 0.000000}
\pgfsetstrokecolor{dialinecolor}
\node[anchor=west] at (47.450000\du,10.325000\du){environnant};
\definecolor{dialinecolor}{rgb}{1.000000, 1.000000, 1.000000}
\pgfsetfillcolor{dialinecolor}
\fill (47.300000\du,10.950000\du)--(47.300000\du,11.950000\du)--(58.195000\du,11.950000\du)--(58.195000\du,10.950000\du)--cycle;
\definecolor{dialinecolor}{rgb}{0.000000, 0.000000, 0.000000}
\pgfsetstrokecolor{dialinecolor}
\draw (47.300000\du,10.950000\du)--(47.300000\du,11.950000\du)--(58.195000\du,11.950000\du)--(58.195000\du,10.950000\du)--cycle;
% setfont left to latex
\definecolor{dialinecolor}{rgb}{0.000000, 0.000000, 0.000000}
\pgfsetstrokecolor{dialinecolor}
\node[anchor=west] at (47.450000\du,11.650000\du){+calculForce(aForme:CForme)};
\pgfsetlinewidth{0.100000\du}
\pgfsetdash{}{0pt}
\definecolor{dialinecolor}{rgb}{1.000000, 1.000000, 1.000000}
\pgfsetfillcolor{dialinecolor}
\fill (13.350100\du,13.237500\du)--(13.350100\du,14.637500\du)--(24.245100\du,14.637500\du)--(24.245100\du,13.237500\du)--cycle;
\definecolor{dialinecolor}{rgb}{0.000000, 0.000000, 0.000000}
\pgfsetstrokecolor{dialinecolor}
\draw (13.350100\du,13.237500\du)--(13.350100\du,14.637500\du)--(24.245100\du,14.637500\du)--(24.245100\du,13.237500\du)--cycle;
% setfont left to latex
\definecolor{dialinecolor}{rgb}{0.000000, 0.000000, 0.000000}
\pgfsetstrokecolor{dialinecolor}
\node at (18.797600\du,14.187500\du){CVent};
\definecolor{dialinecolor}{rgb}{1.000000, 1.000000, 1.000000}
\pgfsetfillcolor{dialinecolor}
\fill (13.350100\du,14.637500\du)--(13.350100\du,15.037500\du)--(24.245100\du,15.037500\du)--(24.245100\du,14.637500\du)--cycle;
\definecolor{dialinecolor}{rgb}{0.000000, 0.000000, 0.000000}
\pgfsetstrokecolor{dialinecolor}
\draw (13.350100\du,14.637500\du)--(13.350100\du,15.037500\du)--(24.245100\du,15.037500\du)--(24.245100\du,14.637500\du)--cycle;
\definecolor{dialinecolor}{rgb}{1.000000, 1.000000, 1.000000}
\pgfsetfillcolor{dialinecolor}
\fill (13.350100\du,15.037500\du)--(13.350100\du,16.037500\du)--(24.245100\du,16.037500\du)--(24.245100\du,15.037500\du)--cycle;
\definecolor{dialinecolor}{rgb}{0.000000, 0.000000, 0.000000}
\pgfsetstrokecolor{dialinecolor}
\draw (13.350100\du,15.037500\du)--(13.350100\du,16.037500\du)--(24.245100\du,16.037500\du)--(24.245100\du,15.037500\du)--cycle;
% setfont left to latex
\definecolor{dialinecolor}{rgb}{0.000000, 0.000000, 0.000000}
\pgfsetstrokecolor{dialinecolor}
\node[anchor=west] at (13.500100\du,15.737500\du){+calculForce(aForme:CForme)};
\pgfsetlinewidth{0.100000\du}
\pgfsetdash{}{0pt}
\definecolor{dialinecolor}{rgb}{1.000000, 1.000000, 1.000000}
\pgfsetfillcolor{dialinecolor}
\fill (41.050100\du,19.087500\du)--(41.050100\du,24.687500\du)--(55.025100\du,24.687500\du)--(55.025100\du,19.087500\du)--cycle;
\definecolor{dialinecolor}{rgb}{0.000000, 0.000000, 0.000000}
\pgfsetstrokecolor{dialinecolor}
\draw (41.050100\du,19.087500\du)--(41.050100\du,24.687500\du)--(55.025100\du,24.687500\du)--(55.025100\du,19.087500\du)--cycle;
% setfont left to latex
\definecolor{dialinecolor}{rgb}{0.000000, 0.000000, 0.000000}
\pgfsetstrokecolor{dialinecolor}
\node at (48.037600\du,20.037500\du){CResultante};
% setfont left to latex
\definecolor{dialinecolor}{rgb}{0.000000, 0.000000, 0.000000}
\pgfsetstrokecolor{dialinecolor}
\node at (48.037600\du,20.812500\du){Somme des};
\definecolor{dialinecolor}{rgb}{0.000000, 0.000000, 0.000000}
\pgfsetstrokecolor{dialinecolor}
\node at (48.037600\du,21.512500\du){torseurs des};
\definecolor{dialinecolor}{rgb}{0.000000, 0.000000, 0.000000}
\pgfsetstrokecolor{dialinecolor}
\node at (48.037600\du,22.212500\du){efforts};
\definecolor{dialinecolor}{rgb}{0.000000, 0.000000, 0.000000}
\pgfsetstrokecolor{dialinecolor}
\node at (48.037600\du,22.912500\du){appliqués au};
\definecolor{dialinecolor}{rgb}{0.000000, 0.000000, 0.000000}
\pgfsetstrokecolor{dialinecolor}
\node at (48.037600\du,23.612500\du){solide en};
\definecolor{dialinecolor}{rgb}{0.000000, 0.000000, 0.000000}
\pgfsetstrokecolor{dialinecolor}
\node at (48.037600\du,24.312500\du){mouvement};
\definecolor{dialinecolor}{rgb}{1.000000, 1.000000, 1.000000}
\pgfsetfillcolor{dialinecolor}
\fill (41.050100\du,24.687500\du)--(41.050100\du,26.487500\du)--(55.025100\du,26.487500\du)--(55.025100\du,24.687500\du)--cycle;
\definecolor{dialinecolor}{rgb}{0.000000, 0.000000, 0.000000}
\pgfsetstrokecolor{dialinecolor}
\draw (41.050100\du,24.687500\du)--(41.050100\du,26.487500\du)--(55.025100\du,26.487500\du)--(55.025100\du,24.687500\du)--cycle;
% setfont left to latex
\definecolor{dialinecolor}{rgb}{0.000000, 0.000000, 0.000000}
\pgfsetstrokecolor{dialinecolor}
\node[anchor=west] at (41.200100\du,25.387500\du){\#fTableau: array \ensuremath{[}1..4\ensuremath{]} of CTorseur};
% setfont left to latex
\definecolor{dialinecolor}{rgb}{0.000000, 0.000000, 0.000000}
\pgfsetstrokecolor{dialinecolor}
\node[anchor=west] at (41.200100\du,26.187500\du){\#fNombreForces: integer};
\definecolor{dialinecolor}{rgb}{1.000000, 1.000000, 1.000000}
\pgfsetfillcolor{dialinecolor}
\fill (41.050100\du,26.487500\du)--(41.050100\du,27.487500\du)--(55.025100\du,27.487500\du)--(55.025100\du,26.487500\du)--cycle;
\definecolor{dialinecolor}{rgb}{0.000000, 0.000000, 0.000000}
\pgfsetstrokecolor{dialinecolor}
\draw (41.050100\du,26.487500\du)--(41.050100\du,27.487500\du)--(55.025100\du,27.487500\du)--(55.025100\du,26.487500\du)--cycle;
% setfont left to latex
\definecolor{dialinecolor}{rgb}{0.000000, 0.000000, 0.000000}
\pgfsetstrokecolor{dialinecolor}
\node[anchor=west] at (41.200100\du,27.187500\du){+somme()};
\pgfsetlinewidth{0.100000\du}
\pgfsetdash{{1.000000\du}{1.000000\du}}{0\du}
\pgfsetdash{{0.400000\du}{0.400000\du}}{0\du}
\pgfsetmiterjoin
\pgfsetbuttcap
{
\definecolor{dialinecolor}{rgb}{0.000000, 0.000000, 0.000000}
\pgfsetfillcolor{dialinecolor}
% was here!!!
\pgfsetarrowsend{to}
\definecolor{dialinecolor}{rgb}{0.000000, 0.000000, 0.000000}
\pgfsetstrokecolor{dialinecolor}
\draw (36.997500\du,3.600000\du)--(36.997500\du,7.350000\du)--(26.115000\du,7.350000\du);
}
% setfont left to latex
\pgfsetlinewidth{0.100000\du}
\pgfsetdash{{0.400000\du}{0.400000\du}}{0\du}
\pgfsetdash{{0.400000\du}{0.400000\du}}{0\du}
\pgfsetmiterjoin
\pgfsetbuttcap
{
\definecolor{dialinecolor}{rgb}{0.000000, 0.000000, 0.000000}
\pgfsetfillcolor{dialinecolor}
% was here!!!
\pgfsetarrowsend{to}
\definecolor{dialinecolor}{rgb}{0.000000, 0.000000, 0.000000}
\pgfsetstrokecolor{dialinecolor}
\draw (36.997500\du,3.600000\du)--(36.997500\du,13.937500\du)--(24.245100\du,13.937500\du);
}
% setfont left to latex
\pgfsetlinewidth{0.100000\du}
\pgfsetdash{{0.400000\du}{0.400000\du}}{0\du}
\pgfsetdash{{0.400000\du}{0.400000\du}}{0\du}
\pgfsetmiterjoin
\pgfsetbuttcap
{
\definecolor{dialinecolor}{rgb}{0.000000, 0.000000, 0.000000}
\pgfsetfillcolor{dialinecolor}
% was here!!!
\pgfsetarrowsend{to}
\definecolor{dialinecolor}{rgb}{0.000000, 0.000000, 0.000000}
\pgfsetstrokecolor{dialinecolor}
\draw (36.997500\du,2.950000\du)--(36.997500\du,7.325000\du)--(47.249700\du,7.325000\du);
}
% setfont left to latex
\pgfsetlinewidth{0.100000\du}
\pgfsetdash{{0.400000\du}{0.400000\du}}{0\du}
\pgfsetdash{{0.400000\du}{0.400000\du}}{0\du}
\pgfsetmiterjoin
\pgfsetbuttcap
{
\definecolor{dialinecolor}{rgb}{0.000000, 0.000000, 0.000000}
\pgfsetfillcolor{dialinecolor}
% was here!!!
\pgfsetarrowsend{to}
\definecolor{dialinecolor}{rgb}{0.000000, 0.000000, 0.000000}
\pgfsetstrokecolor{dialinecolor}
\draw (36.997500\du,3.600000\du)--(36.997500\du,13.937500\du)--(48.037600\du,13.937500\du)--(48.037600\du,19.039368\du);
}
% setfont left to latex
\pgfsetlinewidth{0.100000\du}
\pgfsetdash{{0.400000\du}{0.400000\du}}{0\du}
\pgfsetdash{{0.400000\du}{0.400000\du}}{0\du}
\pgfsetmiterjoin
\pgfsetbuttcap
{
\definecolor{dialinecolor}{rgb}{0.000000, 0.000000, 0.000000}
\pgfsetfillcolor{dialinecolor}
% was here!!!
\pgfsetarrowsend{to}
\definecolor{dialinecolor}{rgb}{0.000000, 0.000000, 0.000000}
\pgfsetstrokecolor{dialinecolor}
\draw (36.997500\du,3.600000\du)--(36.997500\du,16.287500\du)--(28.312500\du,16.287500\du)--(28.312500\du,19.049793\du);
}
% setfont left to latex
\pgfsetlinewidth{0.100000\du}
\pgfsetdash{}{0pt}
\definecolor{dialinecolor}{rgb}{1.000000, 1.000000, 1.000000}
\pgfsetfillcolor{dialinecolor}
\fill (10.100200\du,34.262500\du)--(10.100200\du,37.062500\du)--(23.305200\du,37.062500\du)--(23.305200\du,34.262500\du)--cycle;
\definecolor{dialinecolor}{rgb}{0.000000, 0.000000, 0.000000}
\pgfsetstrokecolor{dialinecolor}
\draw (10.100200\du,34.262500\du)--(10.100200\du,37.062500\du)--(23.305200\du,37.062500\du)--(23.305200\du,34.262500\du)--cycle;
% setfont left to latex
\definecolor{dialinecolor}{rgb}{0.000000, 0.000000, 0.000000}
\pgfsetstrokecolor{dialinecolor}
\node at (16.702700\du,35.212500\du){CSolideMouvement};
% setfont left to latex
\definecolor{dialinecolor}{rgb}{0.000000, 0.000000, 0.000000}
\pgfsetstrokecolor{dialinecolor}
\node at (16.702700\du,35.987500\du){Solide en};
\definecolor{dialinecolor}{rgb}{0.000000, 0.000000, 0.000000}
\pgfsetstrokecolor{dialinecolor}
\node at (16.702700\du,36.687500\du){mouvement};
\definecolor{dialinecolor}{rgb}{1.000000, 1.000000, 1.000000}
\pgfsetfillcolor{dialinecolor}
\fill (10.100200\du,37.062500\du)--(10.100200\du,40.462500\du)--(23.305200\du,40.462500\du)--(23.305200\du,37.062500\du)--cycle;
\definecolor{dialinecolor}{rgb}{0.000000, 0.000000, 0.000000}
\pgfsetstrokecolor{dialinecolor}
\draw (10.100200\du,37.062500\du)--(10.100200\du,40.462500\du)--(23.305200\du,40.462500\du)--(23.305200\du,37.062500\du)--cycle;
% setfont left to latex
\definecolor{dialinecolor}{rgb}{0.000000, 0.000000, 0.000000}
\pgfsetstrokecolor{dialinecolor}
\node[anchor=west] at (10.250200\du,37.762500\du){\#fResultante: CResultante};
% setfont left to latex
\definecolor{dialinecolor}{rgb}{0.000000, 0.000000, 0.000000}
\pgfsetstrokecolor{dialinecolor}
\node[anchor=west] at (10.250200\du,38.562500\du){\#fPositionSolide: CPositionSolide};
% setfont left to latex
\definecolor{dialinecolor}{rgb}{0.000000, 0.000000, 0.000000}
\pgfsetstrokecolor{dialinecolor}
\node[anchor=west] at (10.250200\du,39.362500\du){\#fVitesse: CVitesse};
% setfont left to latex
\definecolor{dialinecolor}{rgb}{0.000000, 0.000000, 0.000000}
\pgfsetstrokecolor{dialinecolor}
\node[anchor=west] at (10.250200\du,40.162500\du){\#fForme: CForme};
\definecolor{dialinecolor}{rgb}{1.000000, 1.000000, 1.000000}
\pgfsetfillcolor{dialinecolor}
\fill (10.100200\du,40.462500\du)--(10.100200\du,41.462500\du)--(23.305200\du,41.462500\du)--(23.305200\du,40.462500\du)--cycle;
\definecolor{dialinecolor}{rgb}{0.000000, 0.000000, 0.000000}
\pgfsetstrokecolor{dialinecolor}
\draw (10.100200\du,40.462500\du)--(10.100200\du,41.462500\du)--(23.305200\du,41.462500\du)--(23.305200\du,40.462500\du)--cycle;
% setfont left to latex
\definecolor{dialinecolor}{rgb}{0.000000, 0.000000, 0.000000}
\pgfsetstrokecolor{dialinecolor}
\node[anchor=west] at (10.250200\du,41.162500\du){+calculPosition()};
\pgfsetlinewidth{0.100000\du}
\pgfsetdash{}{0pt}
\definecolor{dialinecolor}{rgb}{1.000000, 1.000000, 1.000000}
\pgfsetfillcolor{dialinecolor}
\fill (26.800200\du,34.212500\du)--(26.800200\du,35.612500\du)--(33.460200\du,35.612500\du)--(33.460200\du,34.212500\du)--cycle;
\definecolor{dialinecolor}{rgb}{0.000000, 0.000000, 0.000000}
\pgfsetstrokecolor{dialinecolor}
\draw (26.800200\du,34.212500\du)--(26.800200\du,35.612500\du)--(33.460200\du,35.612500\du)--(33.460200\du,34.212500\du)--cycle;
% setfont left to latex
\definecolor{dialinecolor}{rgb}{0.000000, 0.000000, 0.000000}
\pgfsetstrokecolor{dialinecolor}
\node at (30.130200\du,35.162500\du){CMateriau};
\definecolor{dialinecolor}{rgb}{1.000000, 1.000000, 1.000000}
\pgfsetfillcolor{dialinecolor}
\fill (26.800200\du,35.612500\du)--(26.800200\du,37.412500\du)--(33.460200\du,37.412500\du)--(33.460200\du,35.612500\du)--cycle;
\definecolor{dialinecolor}{rgb}{0.000000, 0.000000, 0.000000}
\pgfsetstrokecolor{dialinecolor}
\draw (26.800200\du,35.612500\du)--(26.800200\du,37.412500\du)--(33.460200\du,37.412500\du)--(33.460200\du,35.612500\du)--cycle;
% setfont left to latex
\definecolor{dialinecolor}{rgb}{0.000000, 0.000000, 0.000000}
\pgfsetstrokecolor{dialinecolor}
\node[anchor=west] at (26.950200\du,36.312500\du){\#fIndex: integer};
% setfont left to latex
\definecolor{dialinecolor}{rgb}{0.000000, 0.000000, 0.000000}
\pgfsetstrokecolor{dialinecolor}
\node[anchor=west] at (26.950200\du,37.112500\du){\#fCoefChoc: real};
\definecolor{dialinecolor}{rgb}{1.000000, 1.000000, 1.000000}
\pgfsetfillcolor{dialinecolor}
\fill (26.800200\du,37.412500\du)--(26.800200\du,37.812500\du)--(33.460200\du,37.812500\du)--(33.460200\du,37.412500\du)--cycle;
\definecolor{dialinecolor}{rgb}{0.000000, 0.000000, 0.000000}
\pgfsetstrokecolor{dialinecolor}
\draw (26.800200\du,37.412500\du)--(26.800200\du,37.812500\du)--(33.460200\du,37.812500\du)--(33.460200\du,37.412500\du)--cycle;
\pgfsetlinewidth{0.100000\du}
\pgfsetdash{}{0pt}
\definecolor{dialinecolor}{rgb}{1.000000, 1.000000, 1.000000}
\pgfsetfillcolor{dialinecolor}
\fill (36.950200\du,34.212500\du)--(36.950200\du,35.612500\du)--(42.455200\du,35.612500\du)--(42.455200\du,34.212500\du)--cycle;
\definecolor{dialinecolor}{rgb}{0.000000, 0.000000, 0.000000}
\pgfsetstrokecolor{dialinecolor}
\draw (36.950200\du,34.212500\du)--(36.950200\du,35.612500\du)--(42.455200\du,35.612500\du)--(42.455200\du,34.212500\du)--cycle;
% setfont left to latex
\definecolor{dialinecolor}{rgb}{0.000000, 0.000000, 0.000000}
\pgfsetstrokecolor{dialinecolor}
\node at (39.702700\du,35.162500\du){CVitesse};
\definecolor{dialinecolor}{rgb}{1.000000, 1.000000, 1.000000}
\pgfsetfillcolor{dialinecolor}
\fill (36.950200\du,35.612500\du)--(36.950200\du,38.212500\du)--(42.455200\du,38.212500\du)--(42.455200\du,35.612500\du)--cycle;
\definecolor{dialinecolor}{rgb}{0.000000, 0.000000, 0.000000}
\pgfsetstrokecolor{dialinecolor}
\draw (36.950200\du,35.612500\du)--(36.950200\du,38.212500\du)--(42.455200\du,38.212500\du)--(42.455200\du,35.612500\du)--cycle;
% setfont left to latex
\definecolor{dialinecolor}{rgb}{0.000000, 0.000000, 0.000000}
\pgfsetstrokecolor{dialinecolor}
\node[anchor=west] at (37.100200\du,36.312500\du){\#fX: real};
% setfont left to latex
\definecolor{dialinecolor}{rgb}{0.000000, 0.000000, 0.000000}
\pgfsetstrokecolor{dialinecolor}
\node[anchor=west] at (37.100200\du,37.112500\du){\#fY: real};
% setfont left to latex
\definecolor{dialinecolor}{rgb}{0.000000, 0.000000, 0.000000}
\pgfsetstrokecolor{dialinecolor}
\node[anchor=west] at (37.100200\du,37.912500\du){\#fOmega: real};
\definecolor{dialinecolor}{rgb}{1.000000, 1.000000, 1.000000}
\pgfsetfillcolor{dialinecolor}
\fill (36.950200\du,38.212500\du)--(36.950200\du,38.612500\du)--(42.455200\du,38.612500\du)--(42.455200\du,38.212500\du)--cycle;
\definecolor{dialinecolor}{rgb}{0.000000, 0.000000, 0.000000}
\pgfsetstrokecolor{dialinecolor}
\draw (36.950200\du,38.212500\du)--(36.950200\du,38.612500\du)--(42.455200\du,38.612500\du)--(42.455200\du,38.212500\du)--cycle;
\pgfsetlinewidth{0.100000\du}
\pgfsetdash{}{0pt}
\definecolor{dialinecolor}{rgb}{1.000000, 1.000000, 1.000000}
\pgfsetfillcolor{dialinecolor}
\fill (46.500200\du,34.312500\du)--(46.500200\du,35.712500\du)--(55.855200\du,35.712500\du)--(55.855200\du,34.312500\du)--cycle;
\definecolor{dialinecolor}{rgb}{0.000000, 0.000000, 0.000000}
\pgfsetstrokecolor{dialinecolor}
\draw (46.500200\du,34.312500\du)--(46.500200\du,35.712500\du)--(55.855200\du,35.712500\du)--(55.855200\du,34.312500\du)--cycle;
% setfont left to latex
\definecolor{dialinecolor}{rgb}{0.000000, 0.000000, 0.000000}
\pgfsetstrokecolor{dialinecolor}
\node at (51.177700\du,35.262500\du){CPosition};
\definecolor{dialinecolor}{rgb}{1.000000, 1.000000, 1.000000}
\pgfsetfillcolor{dialinecolor}
\fill (46.500200\du,35.712500\du)--(46.500200\du,39.112500\du)--(55.855200\du,39.112500\du)--(55.855200\du,35.712500\du)--cycle;
\definecolor{dialinecolor}{rgb}{0.000000, 0.000000, 0.000000}
\pgfsetstrokecolor{dialinecolor}
\draw (46.500200\du,35.712500\du)--(46.500200\du,39.112500\du)--(55.855200\du,39.112500\du)--(55.855200\du,35.712500\du)--cycle;
% setfont left to latex
\definecolor{dialinecolor}{rgb}{0.000000, 0.000000, 0.000000}
\pgfsetstrokecolor{dialinecolor}
\node[anchor=west] at (46.650200\du,36.412500\du){\#fXpixel: integer};
% setfont left to latex
\definecolor{dialinecolor}{rgb}{0.000000, 0.000000, 0.000000}
\pgfsetstrokecolor{dialinecolor}
\node[anchor=west] at (46.650200\du,37.212500\du){\#fYpixel: integer};
% setfont left to latex
\definecolor{dialinecolor}{rgb}{0.000000, 0.000000, 0.000000}
\pgfsetstrokecolor{dialinecolor}
\node[anchor=west] at (46.650200\du,38.012500\du){\#fXmetre: real};
% setfont left to latex
\definecolor{dialinecolor}{rgb}{0.000000, 0.000000, 0.000000}
\pgfsetstrokecolor{dialinecolor}
\node[anchor=west] at (46.650200\du,38.812500\du){\#fYmetre: real};
\definecolor{dialinecolor}{rgb}{1.000000, 1.000000, 1.000000}
\pgfsetfillcolor{dialinecolor}
\fill (46.500200\du,39.112500\du)--(46.500200\du,40.112500\du)--(55.855200\du,40.112500\du)--(55.855200\du,39.112500\du)--cycle;
\definecolor{dialinecolor}{rgb}{0.000000, 0.000000, 0.000000}
\pgfsetstrokecolor{dialinecolor}
\draw (46.500200\du,39.112500\du)--(46.500200\du,40.112500\du)--(55.855200\du,40.112500\du)--(55.855200\du,39.112500\du)--cycle;
% setfont left to latex
\definecolor{dialinecolor}{rgb}{0.000000, 0.000000, 0.000000}
\pgfsetstrokecolor{dialinecolor}
\node[anchor=west] at (46.650200\du,39.812500\du){+fonctions de passage()};
\pgfsetlinewidth{0.100000\du}
\pgfsetdash{}{0pt}
\definecolor{dialinecolor}{rgb}{1.000000, 1.000000, 1.000000}
\pgfsetfillcolor{dialinecolor}
\fill (47.150200\du,44.662500\du)--(47.150200\du,46.062500\du)--(54.692700\du,46.062500\du)--(54.692700\du,44.662500\du)--cycle;
\definecolor{dialinecolor}{rgb}{0.000000, 0.000000, 0.000000}
\pgfsetstrokecolor{dialinecolor}
\draw (47.150200\du,44.662500\du)--(47.150200\du,46.062500\du)--(54.692700\du,46.062500\du)--(54.692700\du,44.662500\du)--cycle;
% setfont left to latex
\definecolor{dialinecolor}{rgb}{0.000000, 0.000000, 0.000000}
\pgfsetstrokecolor{dialinecolor}
\node at (50.921450\du,45.612500\du){CPositionSolide};
\definecolor{dialinecolor}{rgb}{1.000000, 1.000000, 1.000000}
\pgfsetfillcolor{dialinecolor}
\fill (47.150200\du,46.062500\du)--(47.150200\du,47.062500\du)--(54.692700\du,47.062500\du)--(54.692700\du,46.062500\du)--cycle;
\definecolor{dialinecolor}{rgb}{0.000000, 0.000000, 0.000000}
\pgfsetstrokecolor{dialinecolor}
\draw (47.150200\du,46.062500\du)--(47.150200\du,47.062500\du)--(54.692700\du,47.062500\du)--(54.692700\du,46.062500\du)--cycle;
% setfont left to latex
\definecolor{dialinecolor}{rgb}{0.000000, 0.000000, 0.000000}
\pgfsetstrokecolor{dialinecolor}
\node[anchor=west] at (47.300200\du,46.762500\du){\#fAngle: real};
\definecolor{dialinecolor}{rgb}{1.000000, 1.000000, 1.000000}
\pgfsetfillcolor{dialinecolor}
\fill (47.150200\du,47.062500\du)--(47.150200\du,47.462500\du)--(54.692700\du,47.462500\du)--(54.692700\du,47.062500\du)--cycle;
\definecolor{dialinecolor}{rgb}{0.000000, 0.000000, 0.000000}
\pgfsetstrokecolor{dialinecolor}
\draw (47.150200\du,47.062500\du)--(47.150200\du,47.462500\du)--(54.692700\du,47.462500\du)--(54.692700\du,47.062500\du)--cycle;
\pgfsetlinewidth{0.100000\du}
\pgfsetdash{{0.400000\du}{0.400000\du}}{0\du}
\pgfsetdash{{0.400000\du}{0.400000\du}}{0\du}
\pgfsetmiterjoin
\pgfsetbuttcap
{
\definecolor{dialinecolor}{rgb}{0.000000, 0.000000, 0.000000}
\pgfsetfillcolor{dialinecolor}
% was here!!!
\pgfsetarrowsend{to}
\definecolor{dialinecolor}{rgb}{0.000000, 0.000000, 0.000000}
\pgfsetstrokecolor{dialinecolor}
\draw (50.900200\du,40.262500\du)--(50.900200\du,42.437323\du)--(50.921450\du,42.437323\du)--(50.921450\du,44.612146\du);
}
% setfont left to latex
\pgfsetlinewidth{0.100000\du}
\pgfsetdash{}{0pt}
\definecolor{dialinecolor}{rgb}{1.000000, 1.000000, 1.000000}
\pgfsetfillcolor{dialinecolor}
\fill (29.100000\du,40.212500\du)--(29.100000\du,41.612500\du)--(39.995000\du,41.612500\du)--(39.995000\du,40.212500\du)--cycle;
\definecolor{dialinecolor}{rgb}{0.000000, 0.000000, 0.000000}
\pgfsetstrokecolor{dialinecolor}
\draw (29.100000\du,40.212500\du)--(29.100000\du,41.612500\du)--(39.995000\du,41.612500\du)--(39.995000\du,40.212500\du)--cycle;
% setfont left to latex
\definecolor{dialinecolor}{rgb}{0.000000, 0.000000, 0.000000}
\pgfsetstrokecolor{dialinecolor}
\node at (34.547500\du,41.162500\du){CForme};
\definecolor{dialinecolor}{rgb}{1.000000, 1.000000, 1.000000}
\pgfsetfillcolor{dialinecolor}
\fill (29.100000\du,41.612500\du)--(29.100000\du,43.412500\du)--(39.995000\du,43.412500\du)--(39.995000\du,41.612500\du)--cycle;
\definecolor{dialinecolor}{rgb}{0.000000, 0.000000, 0.000000}
\pgfsetstrokecolor{dialinecolor}
\draw (29.100000\du,41.612500\du)--(29.100000\du,43.412500\du)--(39.995000\du,43.412500\du)--(39.995000\du,41.612500\du)--cycle;
% setfont left to latex
\definecolor{dialinecolor}{rgb}{0.000000, 0.000000, 0.000000}
\pgfsetstrokecolor{dialinecolor}
\node[anchor=west] at (29.250000\du,42.312500\du){+fMat : array of boolean};
% setfont left to latex
\definecolor{dialinecolor}{rgb}{0.000000, 0.000000, 0.000000}
\pgfsetstrokecolor{dialinecolor}
\node[anchor=west] at (29.250000\du,43.112500\du){+fCentreInertie : CPosition};
\definecolor{dialinecolor}{rgb}{1.000000, 1.000000, 1.000000}
\pgfsetfillcolor{dialinecolor}
\fill (29.100000\du,43.412500\du)--(29.100000\du,44.412500\du)--(39.995000\du,44.412500\du)--(39.995000\du,43.412500\du)--cycle;
\definecolor{dialinecolor}{rgb}{0.000000, 0.000000, 0.000000}
\pgfsetstrokecolor{dialinecolor}
\draw (29.100000\du,43.412500\du)--(29.100000\du,44.412500\du)--(39.995000\du,44.412500\du)--(39.995000\du,43.412500\du)--cycle;
% setfont left to latex
\definecolor{dialinecolor}{rgb}{0.000000, 0.000000, 0.000000}
\pgfsetstrokecolor{dialinecolor}
\node[anchor=west] at (29.250000\du,44.112500\du){+CalcuCentreInertie()};
\end{tikzpicture}

\newpage
\subsection{Impressions écrans}

\begin{center}
\begin{figure}[h]
\begin{center}
\scalebox{0.5}[0.5]{\includegraphics*{../images/global.jpg}}
\end{center}
\caption{Fenêtre principale}
\end{figure}
\end{center}

\begin{center}
\begin{figure}[h]
\begin{center}
\scalebox{0.5}[0.5]{\includegraphics*{../images/dessin_objet.jpg}}
\end{center}
\caption{Fenêtre gérant le dessin de l'objet}
\end{figure}
\end{center}

\begin{center}
\begin{figure}[h]
\begin{center}
\scalebox{0.6}[0.6]{\includegraphics*{../images/dessin_decor.jpg}}
\end{center}
\caption{Fenêtre gérant le dessin du décor}
\end{figure}
\end{center}

\end{document}

